\section{Partially Ordered Datasets (POD)} \label{sec:system}
We define the notion of Partially Ordered Datasets (POD), which allow a set of \emph{clients} to obtain a partially ordered list of transactions from a set of \emph{replicas} with censorship-freeness guarantees and accountability for safety violations. Before we define the syntax and security guarantees for PODs, we present basic notation and the model we consider. 

\subsection{Model}

\dotparagraph{Notation}
We distinguish two types of parties, $n$ \emph{replicas} and an unlimited number of \emph{clients}.
Replicas run the protocol and clients interact with it using its interface.
We call a party -- replica or client -- \emph{correct}, if it follows the protocol, and \emph{malicious} otherwise.

% \begin{remark}[Clients]
% Certain operations are driven by the client.
% However, a malicious client can lose liveness for \emph{its own} operations or transactions, but cannot affect other clients.
% Importantly, a malicious client cannot break any safety property of the system.
% \end{remark}

\dotparagraph{Modeling Time}
We assume that time proceeds in \emph{rounds} and that parties have loosely synchronized clocks allowing them to determine the current round. Each party keeps a local round counter that is used to assign timestamps to transactions and protocol messages as needed. Without loss of generality, our analysis can be extended to a model with loosely synchronized clocks where two honest parties may have two different round counters $r_1,r_2$ such that $|r_1-r_2|=\delta$ for a finite $\delta$.

\dotparagraph{Modeling Network Synchrony}
In the synchronous model, we assume that every message is delivered within a \emph{known} finite delay of $\Delta$ rounds. In the partially synchronous model, we assume that messages are delivered within a known finite delay of $\Delta$ rounds only \emph{after} an unknown but finite period of time called the Global Stabilization Time (GST). Messages sent before the GST may be delivered only $\Delta$ rounds after the GST. Finally, in the asynchronous model, we assume that messages are delivered within an \emph{unknown} but finite delay $\Delta$.

%We use the notion of a \emph{round}: each transaction gets associated with a specific round. We say the transaction becomes \emph{\confirmed} at that round. \bd{There's a confusion between the definition of a round as a unit of time during which communication takes place and the use of rounds as  timestamps for specific transactions. We also lack a definition of synchronous/semi-synchronous communication.}

\subsection{POD Syntax}

% \begin{definition}[Minimum and maximum transaction bounds]\label{def:min-max-bounds}
%     We also associate with each transaction a \emph{minimum round} (\rmin) and a \emph{maximum round} (\rmax).
%     They bound the \confirmed round \rconf for \tx across all clients and are monotonic. Specifically,
%     (1) if an honest client computes some $\rmin$ and $\rmax$ for a transaction
%     and that transaction becomes \confirmed for some (possibly different) client with recorded round $\rconf$, then
%     $\rmin \leq \rconf \leq \rmax$,
%     and (2) if an honest client computes $\rmin$ and $\rmax$, and later updates these values to $\rmin'$ and $\rmax'$,
%     then $\rmin' \geq \rmin$ and $\rmax' \leq \rmax$.
% \end{definition}
% The \rmin, \rmax, \rconf are computed from the votes received for that transaction and they are updated when new votes are received.
% Initially, $\rmin=0$ and $\rmax=\infty$.


\begin{definition}[Transaction and associated values]
    We associate with each transaction \tx values $(\tx, \rmin, \rmax, \rconf, \txcert )$.
    We call $\rmin \in \CN$ the \emph{minimum round}, $\rmax \in \CN \cup \{\infty\}$ the \emph{maximum round},
    and $\rconf \in \CN \cup \{\bot\}$ the \emph{\confirmed round}.
    For the maximum round, symbol~$\infty$ stands for an unbounded large number,
    and for the \confirmed round, symbol~$\bot$ stands for an undefined value.
    \oa{also a filed unconfirmed, \pend, \confirmed.}
    Finally, $\txcert$ is a \emph{transaction certificate} that proves the correctness of $\rmin, \rmax, \rconf$.
    For a \confirmed transaction we have that $\rconf \neq \bot$, and for a transaction that is not \confirmed we have that $\rconf = \bot$.
\end{definition}

\begin{definition}[Partially Ordered Dataset (POD)]
    A \emph{Partially Ordered Dataset (POD)} $\POD$ is a finite sequence of transactions and associated values $(\tx, \rmin, \rmax, \rconf, \txcert )$.
\end{definition}
    
% \begin{definition}[Past-perfect round and past-perfection certificate]
% \label{def:past-perfect}
% We say that a round $\rperf$ becomes \emph{past perfect} in the view of a client
% when the client obtains a past-perfection certificate for \rperf.
% A \emph{past-perfection certificate} \ppcert is a collection of votes from replicas
% that prove that any new transaction that enters the system will be assigned a \confirmed round larger than \rperf.
% Given a past-perfection certificate \ppcert, we refer to the past-perfect round it implies by $\ppcert.\rperf$.
% In the construction, we will have clients also output a POD \POD, containing all transactions they have seen by that point, together with \rperf.
% We refer to that by $\ppcert.\POD$.
% \end{definition}

\begin{definition}[POD protocol]
A POD protocol has the following interface.
\begin{itemize}
    \item $\opwrite{\tx} \returns (\tx, \rmin, \rmax, \rconf, \txcert)$:
    It gets a transaction \tx and outputs values its associated values $\rmin, \rmax$, $\rconf$, and \txcert.

    \item $\opreadall{} \returns \left( \POD, \rperf, \ppcert \right)$:
    It returns a POD \POD, a round number \rperf, called the \emph{past perfect round}, and a \emph{past-perfection certificate} \ppcert that proves the correctness of the \POD and \rperf.

    \item \opidentify{\txcert, \ppcert} \returns $P' \subset P$ 
    
\end{itemize}
\end{definition}

When a client calls $\opreadall{}$ and obtains some \POD and some $(\tx, \rmin, \rmax, \rconf, \txcert) \in \POD$, for which $\rconf \neq \bot$, we say that the client \emph{observes} transaction \tx as \confirmed.
\oa{Same for \pend. }

\subsection{POD Security Guarantees}
A POD satisfies the following properties.

\begin{enumerate}
    \item{Liveness}
    \begin{description}

        \item[Write confirmation:]
        Assume an honest client calls \opwrite{\tx} at round \round.
        Then, the call will return $(\tx, \rmin, \rmax, \rconf, \txcert)$, where $\rconf \neq \bot$.
        Particularly in synchrony, (1) it will return by round $\round + 2\delta$,
        (2) the \confirmed round of \tx will be $\rconf \in [\round, \round + \delta]$,
        (3) the maximum round of \tx will be $\rmax \in [\round, \round + \delta]$,
        (4) it holds that $\rmax - \rmin \leq \delta$.

\bd{Is $\delta$ really the drift between honest party clocks? Then these guarantees for the synchronous are specific to the concrete construction we show, no? It would be better to define this write confirmation property generically and then set these concrete values later when proving the property for the concrete protocol.}


        \item[Transaction liveness:]
        Any transaction written by an honest client eventually becomes \confirmed.
        Particularly in synchrony:
        \begin{enumerate}
            \item The client that writes \tx by calling \opwrite{\tx} at round \round will obtain a certificate \txcert by round $\round + 2\delta$.
            \item Any other client that calls \opread{} at round $\round'$ will obtain \txcert, for any $\round' \geq \round + \delta$.
            \item The certificate \txcert will satisfy $\txcert.\rconf \in [\round, \round + \delta]$
            and $\txcert.\rmax \in [\round, \round + \delta]$ and
            $\rmax - \rmin \leq \delta$.
        \end{enumerate}
        % \item[Transaction liveness:]
        % Assume an honest client calls \opwrite{\tx} at round \round.
        % If an honest client calls \opreadall{} at round $\round'$, then 
        % the response will contain \tx as \confirmed, for some $\round' > \round$.
        % Particularly in synchrony this holds for any $\round' \geq \round + \delta$,
        
        \bd{Same comment about the generic property vs. the concrete parameter in synchrony.}
        
        \item[Past-perfection liveness:]
        In synchrony, rounds become past-perfect after $\delta$ time. That is, if an honest client calls \opreadall{} at round $\round$, then the returned \rperf will satisfy $\rperf \geq \round - \delta$.
    \end{description}

  \bd{Same comment about the generic property vs. the concrete parameter in synchrony.}


    \item{Safety}
    \begin{description}
        \item[Confirmation totality:]
        If an honest client calls \opread{} at round \round and outputs \tx as \confirmed,
        then eventually any other honest client that calls \opread{} at round $\round'$ will output \tx as \confirmed.
        Particularly in synchrony, this holds for any $\round' \geq \round + u$.
        \oa{Proof sketch: a malicious replica returns \tx only to the first client, making the difference
        between being confirmed and not. Using gossip/relaying, after $u$ all honest will have seen \tx, hence the second client also sees it as confirmed. }
        
        \bd{Same comment about the generic property vs. the concrete parameter in synchrony.}
        \oa{If an honest client observes transaction \tx as \confirmed, then eventually any other honest client will.}
        
        \item[Pending-confirmation soundness:]
        \oa{If an honest client observes transaction \tx as \pend, then eventually it will observe it as \confirmed.}

        \item[Bound safety:]
        % Assume a call to \opread{} by an honest client (or the call to \opwrite{} that initially wrote \tx) outputs \tx with some \rmin and \rmax,
        % and a \opread{} by an honest client (or the call to \opwrite{} that initially wrote \tx) outputs \tx with some $\rconf \neq \bot$.
        % Then, the transaction bounds definition (\Cref{def:min-max-bounds}) is satisfied, that is:
        % (1) $\rconf \in [\rmin, \rmax]$, and
        % (2) if the first client updates the bounds to $\rmin'$ and $\rmax'$,then $\rmin' \geq \rmin$ and $\rmax' \leq \rmax$.
        % The transaction bounds \rmin, \rmax satisfy \Cref{def:min-max-bounds}.
        (1) if an honest client computes some $\rmin$ and $\rmax$ for a transaction
        and that transaction becomes \confirmed for some (possibly different) client with recorded round $\rconf$, then
        $\rmin \leq \rconf \leq \rmax$,
        and (2) if an honest client computes $\rmin$ and $\rmax$, and later updates these values to $\rmin'$ and $\rmax'$,
        then $\rmin' \geq \rmin$ and $\rmax' \leq \rmax$.

        \item[Past-perfection safety:]
        Assume a \opreadall{} returns a POD \POD and a past-perfect round \rperf.
        For any transaction $\tx \not \in \POD$, no certificate \txcert for \tx, such that $\txcert.\rconf < \rperf$, can ever be created.

        \item[Fair punishment]
        No honest replica gets punished as a result of malicious operation. If $\opidentify{\txcert, \ppcert} \returns P'$, where \txcert is a record certificate for transaction \tx
        and $\ppcert$ is a past-perfect certificate for a POD \POD,
        can only be created if all parties in $P'$ sign \tx and \POD. \oa{TODO: Update}
    \end{description}
\end{enumerate}

\begin{lemma}
    If an honest client writes \tx in a synchronous period and two honest clients compute \confirmed rounds $\rconf$ and $\rconf'$, respectively,
    then $|\rconf - \rconf'| \leq \delta$.
\end{lemma}
\begin{proof}
    It follows from the \emph{transaction liveness} property.
\end{proof}

  \bd{Same comment about the generic property vs. the concrete parameter in synchrony.}

\bd{We can probably show that these are upper bounds in the synchronous case, but again I think it's better to define all properties generically and then consider the concrete parameters for our construction and work out the bounds for different synchrony regimes separately.}

% \subsection{broadcast-deliver}
% The system exposes the method \opbcast{\tx}, which allows a user to input a transaction \tx to the system, and returns a certificate \txcert.
% The method also allows the caller to specify a set of transactions that have already been delivered, by calling \opbcast{\tx, [$\deptxid_1, \ldots, \deptxid_\ell$]}.
% The system also exposes an event \opdel{\tx}, which is observed by users who monitor the system.
% It satisfies the following properties.
% \begin{description}
%     \item[Censorship-resistance]
%     If a user calls \opbcast{\tx} then eventually (1) that user gets \txcert as output and (2) all honest users output \opdel{\tx}.
% \end{description}


